\chapter{Cockpit décisionnel Knowage BI}
\label{ch:knowage}

\section{Introduction}
Conformément à la dernière étape de la méthodologie GIMSI (mise en œuvre et amélioration permanente), ce chapitre présente le cockpit décisionnel développé avec la plateforme \textbf{Knowage}, outil de Business Intelligence open source connecté directement à l'entrepôt de données PostgreSQL. Ce cockpit restitue, sous forme de tableaux de bord interactifs, les indicateurs spécifiés au chapitre~\ref{ch:analyse}.

\section{Organisation du cockpit}
Le cockpit est structuré en plusieurs feuillets (\emph{sheets}) thématiques, chacun ciblant un axe d'analyse~: accueil et navigation, trafic, déplacements, démographie, et accessibilité.

\subsection{Indicateurs de trafic}
Cette section restitue les volumes de comptage agrégés par site et par tranche horaire, permettant d'identifier les heures de pointe et les sites les plus chargés du réseau.

\begin{figure}[H]
\centering
\includegraphics[width=0.8\textwidth]{kpis/DS_Trafic_Horaire.png}
\caption{Évolution horaire du trafic}
\end{figure}

\begin{figure}[H]
\centering
\includegraphics[width=0.45\textwidth]{kpis/DS_Trafic_ParSite.png}
\hfill
\includegraphics[width=0.45\textwidth]{kpis/DS_Trafic_ParSens.png}
\caption{Volume de trafic par site et par sens de circulation}
\end{figure}

\begin{figure}[H]
\centering
\includegraphics[width=0.7\textwidth]{kpis/DS_HeuresPointe.png}
\caption{Identification des heures de pointe}
\end{figure}

\subsection{Indicateurs de déplacement}
Cette section analyse les déplacements individuels déclarés~: temps de parcours par motif, répartition des flux par zone et coût moyen des déplacements.

\begin{figure}[H]
\centering
\includegraphics[width=0.45\textwidth]{kpis/DS_TempsParcours_Motif.png}
\hfill
\includegraphics[width=0.45\textwidth]{kpis/DS_CoutMoyen.png}
\caption{Temps de parcours par motif et coût moyen des déplacements}
\end{figure}

\begin{figure}[H]
\centering
\includegraphics[width=0.7\textwidth]{kpis/DS_FluxZones.png}
\caption{Flux de déplacements entre zones géographiques}
\end{figure}

\subsection{Indicateurs démographiques}
Cette section dresse le profil sociodémographique des individus enquêtés~: sexe, statut matrimonial, niveau d'instruction, statut d'activité et possession du permis de conduire.

\begin{figure}[H]
\centering
\includegraphics[width=0.45\textwidth]{kpis/DS_Demo_Sexe.png}
\hfill
\includegraphics[width=0.45\textwidth]{kpis/DS_Demo_StatutMatrimonial.png}
\caption{Répartition par sexe et statut matrimonial}
\end{figure}

\begin{figure}[H]
\centering
\includegraphics[width=0.45\textwidth]{kpis/DS_Demo_Instruction.png}
\hfill
\includegraphics[width=0.45\textwidth]{kpis/DS_Demo_Activite.png}
\caption{Niveau d'instruction et statut d'activité}
\end{figure}

\subsection{Indicateurs d'accessibilité}
Cette dernière section restitue les indicateurs d'accessibilité aux services et aux transports en commun, directement issus du score d'accessibilité composite défini au chapitre~\ref{ch:analyse}.

\begin{figure}[H]
\centering
\includegraphics[width=0.45\textwidth]{kpis/DS_Acces_Services.png}
\hfill
\includegraphics[width=0.45\textwidth]{kpis/DS_Acces_ManqueTransport.png}
\caption{Accès aux services et manque de transport déclaré}
\end{figure}

\begin{figure}[H]
\centering
\includegraphics[width=0.45\textwidth]{kpis/DS_Acces_Difficultes.png}
\hfill
\includegraphics[width=0.45\textwidth]{kpis/DS_Acces_Inondation.png}
\caption{Difficultés d'accès déclarées et exposition aux inondations}
\end{figure}

\section{Synthèse des KPIs produits}
Le tableau~\ref{tab:kpis_synth} récapitule les principaux indicateurs clés de performance restitués dans le cockpit Knowage, calculés à partir des volumes réels de l'entrepôt de données.

\begin{table}[H]
\centering
\caption{Synthèse des indicateurs clés restitués (cockpit Knowage)}
\label{tab:kpis_synth}
\begin{tabular}{lr}
\toprule
\textbf{Indicateur} & \textbf{Valeur} \\
\midrule
Individus enquêtés (EMD) & 37\,248 \\
Ménages enquêtés (EMD) & 4\,521 \\
Sites de comptage de trafic & 49 \\
Observations de trafic & 775\,480 \\
Zones géographiques analysées & 507 \\
\bottomrule
\end{tabular}
\end{table}

\section{Rendu du cockpit -- captures}

Les figures suivantes illustrent le rendu réel du cockpit Knowage, avec ses six feuillets thématiques, dans les deux thèmes visuels disponibles (clair et sombre).

\begin{figure}[H]
\centering
\includegraphics[width=0.48\textwidth]{cockpit_accueil_light.png}
\hfill
\includegraphics[width=0.48\textwidth]{cockpit_accueil_dark.png}
\caption{Sheet \emph{Accueil} du cockpit Knowage -- mode clair (gauche) et mode sombre (droite)}
\label{fig:cockpit_accueil}
\end{figure}

\begin{figure}[H]
\centering
\includegraphics[width=0.48\textwidth]{cockpit_trafic_light.png}
\hfill
\includegraphics[width=0.48\textwidth]{cockpit_deplacements.png}
\caption{Sheet \emph{Trafic} (gauche) et sheet \emph{Déplacements} (droite)}
\label{fig:cockpit_trafic_dep}
\end{figure}

\begin{figure}[H]
\centering
\includegraphics[width=0.48\textwidth]{cockpit_demographie.png}
\hfill
\includegraphics[width=0.48\textwidth]{cockpit_accessibilite.png}
\caption{Sheet \emph{Démographie} (gauche) et sheet \emph{Accessibilité} (droite)}
\label{fig:cockpit_demo_acc}
\end{figure}

\begin{figure}[H]
\centering
\includegraphics[width=0.7\textwidth]{cockpit_ia.png}
\caption{Sheet \emph{IA} -- restitution synthétique des résultats de Machine Learning}
\label{fig:cockpit_ia}
\end{figure}

\section{Apport du cockpit pour les décideurs}
Le cockpit Knowage répond directement aux besoins des deux profils de décideurs identifiés~: le \textbf{directeur de planification} y visualise la répartition géographique et temporelle des déplacements pour orienter ses choix stratégiques (politique de lignes, investissements d'infrastructure)~; le \textbf{chef d'exploitation} y identifie les heures de pointe et les anomalies de trafic pour ajuster en temps réel les fréquences de desserte et l'affectation des ressources.

\section{Conclusion}
Le cockpit Knowage constitue le point de restitution central du système décisionnel, transformant les données brutes de l'entrepôt en informations directement exploitables par les deux profils de décideurs du CETUD. Le chapitre suivant présente l'application web développée pour compléter ce dispositif, à destination à la fois des citoyens et des décideurs.

\clearpage
